\documentclass[11pt]{article}
\usepackage{amsmath,amssymb,mathtools}
\usepackage{geometry}
\usepackage{bm}
\geometry{a4paper,margin=2.5cm}
\allowdisplaybreaks

\title{STEP2: Multi-Input Merging and Continued Propagation\\
Arbitrary Fixed-Proportion Input/Output Recipes}
\author{}
\date{}

\begin{document}
\maketitle

\section{Purpose: Merging Does Not End Propagation}

STEP1 has already established that:

\begin{itemize}
    \item A node's own production-start event is described by
    \(\sigma_p,\rho_p\) or the general \(D_p^{act}(t)\);
    \item The supply deficit visible to a specific downstream node from a particular upstream node is described by
    \(D_{h\to p}(t)\);
    \item Production delay, transportation delay, and inventories along the path
    are all handled by branch propagation in STEP1.
\end{itemize}

STEP2 answers only the following question:
\[
\boxed{\begin{gathered}
\text{After multiple required inputs reach the same node,}\\
\text{how do they determine its production-start activity?}
\end{gathered}}
\]

After the local activity is obtained, propagation continues: the result is passed back to STEP1 for propagation to all downstream nodes.

Thus the overall interface is
\[
\boxed{
\text{STEP1 propagation}
\rightarrow
\text{STEP2 merge}
\rightarrow
\text{STEP1 propagation}
\rightarrow\cdots
}
\tag{1}
\]

\section{General Production Node: Arbitrary Fixed-Proportion Input/Output Recipes}

Consider a production node \(P_p\).

Its set of required inputs is
\[
\mathcal I_p=\{1,\ldots,n_p\},
\]
and its set of outputs is
\[
\mathcal O_p=\{1,\ldots,m_p\}.
\]

The fixed recipe is written as
\[
\boxed{
\sum_{h\in\mathcal I_p}a_hX_h
\longrightarrow
\sum_{r\in\mathcal O_p}b_rY_r,
}
\tag{2}
\]
where
\[
a_h>0,
\qquad
b_r>0.
\]

This is not restricted to
\[
1+1\to1,
\]
but allows any finite number of required inputs and any finite number of outputs,
with arbitrary fixed positive input/output coefficients.

Let the initialized normal production intensity be
\[
\boxed{
u_p^0\equiv u_p(0).
}
\]

Then the initialized normal consumption rate for input \(h\) is
\[
\boxed{
f_{p,h}^0=a_hu_p^0,
}
\tag{3}
\]
and the initialized normal completion rate for output \(r\) is
\[
\boxed{
g_{p,r}^0=b_ru_p^0.
}
\tag{4}
\]

We adopt a fixed-proportion / Leontief production recipe:
All required inputs are indispensable,
and neither input substitution nor dynamic recipe switching is permitted at this stage.

\section{Inputs Passed from STEP1 to STEP2}

For each
\[
h\in\mathcal I_p,
\]
suppose required input \(h\) corresponds to a branch whose propagation has reached the input-receiving plane of \(P_p\).

STEP1 returns
\[
\boxed{
D_{h\to p}(t).
}
\tag{5}
\]

Its meaning is:
\[
\boxed{\begin{gathered}
D_{h\to p}(t)=\\
\text{cumulative equivalent normal-production time missing}\\
\text{for input \(h\), as seen by \(P_p\) by time \(t\)}
\end{gathered}}
\]

Therefore the corresponding cumulative material deficit is
\[
\boxed{
a_hu_p^0D_{h\to p}(t).
}
\tag{6}
\]

All
\[
D_{h\to p}(t)
\]
must already refer to the same observation plane:
\[
\boxed{\begin{gathered}
\text{the production-start plane where \(P_p\) can access}\\
\text{all required inputs and decide whether to initiate production}
\end{gathered}}
\]

STEP2 does not recalculate anything already covered by STEP1,
including production delay, transportation delay, or inventory coverage.

\section{Single-Window Input Branches and General Input Branches}

If input branch \(h\) has only one effective supply-interruption window
\[
[t_{h\to p}^{\mathrm{stop}},
t_{h\to p}^{\mathrm{rec}}),
\]
then
\[
\boxed{
D_{h\to p}(t)
=
\min\left\{
(t-t_{h\to p}^{\mathrm{stop}})_+,
\Delta_{h\to p}
\right\},
}
\tag{7}
\]
where
\[
\boxed{
\Delta_{h\to p}
=
\left(
t_{h\to p}^{\mathrm{rec}}
-
t_{h\to p}^{\mathrm{stop}}
\right)_+.
}
\tag{8}
\]

However, in a general DAG an input branch may already contain multiple disjoint deficits,
so the main interface between STEP1 and STEP2 is always the complete function
\[
\boxed{D_{h\to p}(t),}
\]
rather than merely a pair
\[
(t^{\mathrm{stop}},t^{\mathrm{rec}}).
\]

\section{Why Merging Is Not a Simple Union of Windows}

If any required input temporarily prevents \(P_p\) from producing,
other required inputs will also stop being consumed by the production process.

These nonbinding inputs may keep arriving and be stored during the shutdown,
creating additional buffers for later periods.

Thus, shortage windows computed independently for different input branches must not be
simply unioned, nor should their shortage durations be directly added.

The correct objects are the individual inputs'
\[
D_{h\to p}(t),
\]
which are compared as cumulative constraints at the production-start plane.

\section{Main STEP2 Theorem: Upper Envelope of Required-Input Constraints}

By time \(t\), if the production node had operated continuously at its normal intensity,
the cumulative consumption of input \(h\) would have been
\[
a_hu_p^0t.
\]

Because this input branch has a cumulative shortfall of
\[
a_hu_p^0D_{h\to p}(t),
\]
the maximum cumulative production activity that this input can support by time \(t\) is
\[
\frac{
a_hu_p^0t
-
a_hu_p^0D_{h\to p}(t)
}{a_h}.
\]

Hence
\[
\boxed{
\frac{
a_hu_p^0t
-
a_hu_p^0D_{h\to p}(t)
}{a_h}
=
u_p^0
\left[
t-D_{h\to p}(t)
\right].
}
\tag{9}
\]

Let the cumulative production activity initiated by node \(P_p\) up to time \(t\) be
\[
\boxed{
U_p(t)=\int_0^t u_p(s)\,ds.
}
\tag{10}
\]

Since all required inputs must be simultaneously satisfied,
\[
U_p(t)
\le
u_p^0
\left[
t-D_{h\to p}(t)
\right]
\qquad
\forall h\in\mathcal I_p.
\]

Under the current greedy baseline,
the node produces at normal intensity whenever materials permit, so actual cumulative activity takes the tightest constraint:
\[
\begin{aligned}
U_p(t)
&=
\min_{h\in\mathcal I_p}
u_p^0
\left[
t-D_{h\to p}(t)
\right]
\\
&=
u_p^0
\left[
t-
\max_{h\in\mathcal I_p}
D_{h\to p}(t)
\right].
\end{aligned}
\tag{11}
\]

Define the node's own cumulative production-start deficit
\[
\boxed{
D_p^{act}(t)
=
\max_{h\in\mathcal I_p}
D_{h\to p}(t).
}
\tag{12}
\]

Then
\[
\boxed{
U_p(t)
=
u_p^0
\left[
t-D_p^{act}(t)
\right].
}
\tag{13}
\]

Away from the switching points of events,
\[
\boxed{
u_p(t)
=
u_p^0
\left[
1-\dot D_p^{act}(t)
\right].
}
\tag{14}
\]

Under the complete-production / complete-shutdown baseline,
\[
\dot D_p^{act}(t)\in\{0,1\}
\]
holds almost everywhere.

\section{Why Arbitrary Input Coefficients \(a_h\) Do Not Change the \(\max\) Structure}

Equation \((9)\) shows that
each required input's material deficit is divided by its own recipe coefficient \(a_h\),
so that all constraints are expressed in the same units:
\[
\boxed{\text{Equivalent time of normal production activity}.}
\]

Thus although
\[
a_1,a_2,\ldots,a_{n_p}
\]
may be arbitrarily different,
after normalization each required input yields an
\[
t-D_{h\to p}(t)
\]
activity constraint of this form.

Therefore
\[
\boxed{
\text{arbitrary fixed-proportion required-input coefficients}
\quad\Longrightarrow\quad
D_p^{act}(t)=\max_hD_{h\to p}(t).
}
\tag{15}
\]

The input coefficients \(a_h\) still affect normal flow rates,
\[
f_{p,h}^0=a_hu_p^0
\]
and thereby material quantities and inventory coverage,
but not the normalized \(\max\)-merge form.

\section{Arbitrary Output Coefficients and Branch Allocation}

After merging, \(P_p\) has just one common production-start activity
\[
u_p(t),
\qquad
D_p^{act}(t).
\]

The completion rate of output \(r\) is
\[
g_{p,r}(t)=b_ru_p(t).
\]

If that output is then allocated at a fixed proportion
\[
\rho_{p,r\to q}
\]
to downstream node \(P_q\),
the normal flow rate of that output branch is
\[
\boxed{
f_{p,r\to q}^0
=
\rho_{p,r\to q}b_ru_p^0.
}
\tag{16}
\]

Thus different outputs may have arbitrary fixed positive coefficients and different fixed allocation ratios.
These coefficients scale material quantities and enter each branch's inventory coverage,
but all outputs share the same local production activity
\[
D_p^{act}(t).
\]

\section{Reading \(P_p\)'s Own Stop/Recovery Events from the Merged Result}

From
\[
D_p^{act}(t)
\]
read all maximal connected intervals satisfying
\[
\dot D_p^{act}(t)=1
\]
as:
\[
\boxed{
I_p^{act,\ell}
=
[\sigma_p^\ell,\rho_p^\ell),
\qquad
\ell=1,\ldots,L_p.
}
\tag{17}
\]

These intervals represent
\[
\boxed{
\text{production-start gaps when \(P_p\) itself cannot initiate new production}.
}
\]

These
\[
\sigma_p^\ell,\rho_p^\ell
\]
still are not stop/recovery events propagated to an arbitrary downstream node;
the specific downstream node must be processed again through STEP1.

If the merged result has only one internal window
\[
[\sigma_p,\rho_p),
\]
it reduces exactly to the single-window input format of STEP1.

\section{How to Continue Propagating Shocks Downstream After Merging}

Suppose an output of \(P_p\) supplies downstream node \(P_q\) along a direct branch.

STEP2's output is not an endpoint; it is
\[
\boxed{
D_p^{act}(t),
}
\]
and immediately becomes the input to the next STEP1 operation.

Define the output branch's
\[
\Theta_{p\to q}
=
\tau_p^{prod}
+
\sum_{e\in\mathcal E_{p\to q}}
\tau_e^{trans},
\]
and
\[
C_{p\to q}
=
\sum_{m\in\mathcal M_{p\to q}}
\frac{q_m(0)}{f_m^{out,0}}.
\]

Then
\[
\boxed{
D_{p\to q}(t)
=
\mathcal T_{p\to q}[D_p^{act}](t)
=
\left[
D_p^{act}(t-\Theta_{p\to q})
-
C_{p\to q}
\right]_+.
}
\tag{18}
\]

Therefore
\[
\boxed{
\{D_{h\to p}\}_h
\overset{\mathrm{STEP2}}{\longrightarrow}
D_p^{act}
\overset{\mathrm{STEP1}}{\longrightarrow}
D_{p\to q}.
}
\tag{19}
\]

\section{Recovering the First Version's Symmetric Downstream Formulas for a Single Merged Window}

If \(P_p\) has only one production-start gap after merging in STEP2,
\[
[\sigma_p,\rho_p),
\]
and the gap is long enough to penetrate all inventory on the \(P_p\to P_q\) branch, i.e.,
\[
\rho_p-\sigma_p>C_{p\to q},
\]
STEP1 immediately gives
\[
\boxed{
\begin{aligned}
t_{p\to q}^{\mathrm{stop}}
&=
\sigma_p
+\Theta_{p\to q}
+C_{p\to q},
\\[1mm]
t_{p\to q}^{\mathrm{rec}}
&=
\rho_p
+\Theta_{p\to q}.
\end{aligned}
}
\tag{20}
\]

Thus even when \(P_p\) itself merges multiple inputs,
as long as the local result is a single window,
the same formulas still hold exactly for any downstream node:
\[
\boxed{
\text{stop}=\text{local stop}+\text{delay}+\text{inventory},
\qquad
\text{recovery}=\text{local recovery}+\text{delay}.
}
\tag{21}
\]

If
\[
\rho_p-\sigma_p\le C_{p\to q},
\]
then the output branch's inventory fully absorbs the internal shutdown,
and no effective supply interruption occurs at \(P_q\).

\section{Continued Propagation of Multi-Window Merge Results}

If
\[
D_p^{act}(t)
\]
contains multiple disjoint increasing intervals,
one cannot apply the initial inventory \(C_{p\to q}\) separately to each
\[
[\sigma_p^\ell,\rho_p^\ell)
\]
interval.

The correct procedure is to evaluate the entire cumulative function once:
\[
\boxed{
D_{p\to q}(t)
=
[D_p^{act}(t-\Theta_{p\to q})-C_{p\to q}]_+.
}
\tag{22}
\]

Then, at the receiving plane of \(P_q\), read
\[
\boxed{
[t_{p\to q}^{\mathrm{stop},j},
t_{p\to q}^{\mathrm{rec},j})
=
\text{the \(j\)-th maximal connected interval with \(\dot D_{p\to q}(t)=1\)}.
}
\tag{23}
\]

This prevents the same downstream inventory from being erroneously reinitialized for multiple windows.

\section{One Merge Node Can Cause Different Shocks at Different Downstream Nodes}

If \(P_p\) supplies
\[
P_{q_1},P_{q_2},\ldots,P_{q_s},
\]
then these nodes share the same
\[
D_p^{act}(t),
\]
but each output branch has its own
\[
\Theta_{p\to q_j},
\qquad
C_{p\to q_j}.
\]

Therefore
\[
\boxed{
D_{p\to q_j}(t)
=
\left[
D_p^{act}(t-\Theta_{p\to q_j})
-
C_{p\to q_j}
\right]_+,
\qquad
j=1,\ldots,s.
}
\tag{24}
\]

Consequently, the following can occur:
\[
P_p\text{ itself has already shut down},
\]
yet one downstream node with ample inventory has felt no supply interruption,
while another with little inventory has already experienced a shortage.

Therefore
\[
\boxed{
\text{The local activity event is shared by the node;}
\text{the downstream-visible event is branch-specific.}
}
\tag{25}
\]

\section{Two-Input Example: Shutdown, Recovery, and Another Shutdown After Merging}

Consider
\[
A+B\longrightarrow Y.
\]

Suppose STEP1 has already propagated both upstream branches to \(Y\)'s input-receiving plane:
\[
D_{A\to Y}(t)
=
\begin{cases}
0, & t<2,\\
t-2, & 2\le t<5,\\
3, & t\ge5,
\end{cases}
\]
and
\[
D_{B\to Y}(t)
=
\begin{cases}
0, & t<4,\\
t-4, & 4\le t<8,\\
4, & t\ge8.
\end{cases}
\]

STEP2 merge:
\[
\boxed{
D_Y^{act}(t)
=
\max\left\{
D_{A\to Y}(t),
D_{B\to Y}(t)
\right\}.
}
\]

We obtain
\[
D_Y^{act}(t)
=
\begin{cases}
0, & t<2,\\
t-2, & 2\le t<5,\\
3, & 5\le t<7,\\
t-4, & 7\le t<8,\\
4, & t\ge8.
\end{cases}
\]

Thus \(Y\)'s own production-start gaps are
\[
\boxed{
[2,5)
\quad\text{and}\quad
[7,8).
}
\]

To determine how these internal gaps affect some downstream node \(Z\),
we must not stop at STEP2, but must continue with STEP1:
\[
\boxed{
D_{Y\to Z}(t)
=
\left[
D_Y^{act}(t-\Theta_{Y\to Z})
-C_{Y\to Z}
\right]_+.
}
\tag{26}
\]

Only the increasing intervals read from \(D_{Y\to Z}\)
represent the true events observed by \(Z\):
\[
t_{Y\to Z}^{\mathrm{stop},j},
\qquad
t_{Y\to Z}^{\mathrm{rec},j}.
\]

\section{Correspondence with ``Production Status of Plant \(p\) on Day \(k\)''}

The STEP2 result
\[
D_p^{act}(t)
\]
directly determines the production status of \(P_p\):
\[
\boxed{
\frac{u_p(t)}{u_p^0}
=
1-\dot D_p^{act}(t).
}
\tag{27}
\]

If time is discretized by day with no additional intraday switching on day \(k\),
define
\[
\boxed{
x_p[k]
=
1-
\left(
D_p^{act}(k+1)-D_p^{act}(k)
\right).
}
\tag{28}
\]

Under the binary baseline,
\[
x_p[k]\in\{0,1\},
\]
they respectively indicate normal production or complete shutdown on that day.

Thus, STEP2's \(\max\) merge
does not merely give the ``merge result'';
it also directly gives the original question's target after the merge node:
\[
\boxed{\text{Production status at every time / on every day}.}
\]

\section{Unified Recursion for an Arbitrary DAG}

For any production node \(P_p\), the calculation follows this fixed order.

\subsection*{Input Side: Apply STEP1 First to Every Required-Input Branch}

For each
\[
h\in\mathcal I_p,
\]
calculate
\[
\boxed{
D_{h\to p}(t)
=
\mathcal T_{h\to p}[D_h^{act}](t).
}
\tag{29}
\]

\subsection*{At the Node: STEP2 Performs the Required-Input Merge}

\[
\boxed{
D_p^{act}(t)
=
\max_{h\in\mathcal I_p}
D_{h\to p}(t).
}
\tag{30}
\]

And from
\[
\boxed{
u_p(t)
=
u_p^0
[1-\dot D_p^{act}(t)]
}
\tag{31}
\]
obtain the local production status.

\subsection*{Output Side: Apply STEP1 Again to Every Downstream Branch}

For each direct downstream \(P_q\),
\[
\boxed{
D_{p\to q}(t)
=
\mathcal T_{p\to q}[D_p^{act}](t).
}
\tag{32}
\]

Thus the entire DAG repeatedly uses just two operations:
\[
\boxed{
\text{Along a branch: delay shift + inventory absorption},
}
\]
and
\[
\boxed{
\text{At a production node: required-input \(\max\) merge}.
}
\]

This can be written as
\[
\boxed{
\mathcal T
\rightarrow
\max
\rightarrow
\mathcal T
\rightarrow
\max
\rightarrow\cdots.
}
\tag{33}
\]

\section{Consistency with Single-Input Nodes}

If
\[
|\mathcal I_p|=1,
\]
then
\[
\boxed{
D_p^{act}(t)
=
D_{h\to p}(t).
}
\tag{34}
\]

Thus STEP2 automatically reduces to the identity mapping.
The entire network then requires only STEP1 operator composition along linear chains,
reducing to the original segment-by-segment cumulative formulas.

\section{Inventory Ownership and Avoiding Double Counting}

We must enforce
\[
\boxed{
\text{One physical inventory belongs to exactly one propagation segment}.
}
\]

If an inventory has already been included in an input branch's
\[
C_{h\to p},
\]
STEP2 must not subtract it again as a new buffer.

If \(P_p\) has separate input-side local storage,
it should be represented explicitly as an inventory node \(M\) in the topology,
with clear ownership by a particular STEP1 propagation segment.

Similarly,
if finished-goods storage absorbs an output deficit of \(P_p\),
it should be handled as an inventory node in output-side STEP1,
not hidden within production delay or the STEP2 merge.

\section{Applicability Limits of the Current Closed-Form \(\max\) Rule}

The equation
\[
D_p^{act}(t)=\max_hD_{h\to p}(t)
\]
relies on the following baseline:
\begin{itemize}
    \item All inputs are fixed-proportion required inputs;
    \item Inputs cannot directly substitute for one another;
    \item Production recipe coefficients \(a_h,b_r\) are fixed;
    \item Normally, the node operates at \(u_p^0\) whenever materials permit;
    \item If one input causes a joint shutdown, other arriving inputs can be stored;
    \item Storage-capacity limits, losses, scrapping, and shelf-life constraints are not yet considered;
    \item No active restocking, overproduction, or rerouting;
    \item The current focus is complete shutdown / complete recovery.
\end{itemize}

Thus ``arbitrary input/output ratios'' here means precisely
\[
\boxed{
\text{a multi-input / multi-output Leontief recipe with fixed positive coefficients},
}
\]
not an arbitrary substitutable, switchable, or nonlinear production technology.

If nonbinding inputs cannot continue to be stored during a joint shutdown,
for example, if local storage capacity is binding,
then the simple
\[
D_p^{act}(t)=\max_hD_{h\to p}(t)
\]
may fail, and explicit inventory-conservation equations are needed.

If partial shortages are allowed,
the rule can be generalized to
\[
0\le\dot D(t)\le1,
\]
In that case
\[
u_p(t)
=
u_p^0
[1-\dot D_p^{act}(t)]
\]
naturally represents partial production.

\section{Final STEP2 Interface}

STEP2 can be summarized as
\[
\boxed{
\underbrace{
\{D_{h\to p}(t)\}_{h\in\mathcal I_p}
}_{\text{required-input shortages propagated to \(P_p\) by STEP1}}
\quad
\xrightarrow{\ \max\ }
\quad
\underbrace{
D_p^{act}(t)
}_{\text{local production-start activity gap of \(P_p\)}}
}
\tag{35}
\]

where
\[
\boxed{
D_p^{act}(t)
=
\max_{h\in\mathcal I_p}
D_{h\to p}(t).
}
\]

Then for every direct downstream \(P_q\):
\[
\boxed{
D_p^{act}(t)
\quad
\xrightarrow{\ \mathrm{STEP1}\ }
\quad
D_{p\to q}(t)
=
[D_p^{act}(t-\Theta_{p\to q})-C_{p\to q}]_+.
}
\tag{36}
\]

The final closed loop is therefore
\[
\boxed{\begin{gathered}
\text{upstream local activity}
\rightarrow
\text{STEP1 propagation}
\rightarrow
\text{required-input shortages}\\
\rightarrow
\text{STEP2 max merge}
\rightarrow
\text{local production activity}
\rightarrow
\text{STEP1 propagation}
\end{gathered}}
\]

For a single window, one can always return to the first version's symmetric formulas;
For multiple windows, preserve the entire \(D(t)\) function for continued recursion.

\end{document}
